% !TeX program = xelatex
\documentclass{UAmathtalk}
\usepackage{setspace}
\renewcommand{\when}{12:00 p.m.}
\renewcommand{\where}{Hudson 0110}
%\renewcommand{\where}{Hudson 0110.\\ \href{https://albany.zoom.us/j/81219478289?pwd=N5fy7vxKBXqlQ4mAq4KaeaAM5iAHi4.1}{Zoom ID: 812\,1947\,8289 - Passcode: 233444}}
\author{Marcus Paone, Justin Curry}
\address{SUNY Albany}
\urladdr{https://www.website.com/}
\title{“Perceptrons” Revisited: Learning Homology and the Euler Characteristic
}
\date{Friday, October 2nd, 2026}

\begin{document}

\maketitle
%\renewcommand*{\abstractsize}{\normalsize}
\renewcommand*{\abstractsize}{\fontsize{13.5}{16.2}\selectfont}
\begin{abstract}
    {
        Minsky and Papert’s landmark 1969 text “Perceptrons” is often blamed for causing the first AI winter because it showed that a single layer neural net can only learn the Euler characteristic of a subset of the plane.
        In this talk we will review some of the impossibility theorems for learning topological features and then transition to testing these ideas computationally.
        This will involve a brief detour in random topology for binary image data and then a comparison of what a 2-layer MLP can learn.
        We’ll show results for predicting---without using homology---$\beta_0$, $\beta_1$, and $\chi$ for all $5\times 5$ grid pixel images.
        If time permits, we will discuss whether neural nets can learn Euler’s formula from the latent representations of these images alone.
    }
\end{abstract}
\end{document}


